\documentclass{ximera}



% MATH -----------------------------------------------------------
\newcommand{\nats}{\mathbb N}
\newcommand{\ints}{\mathbb Z}
\newcommand{\rats}{\mathbb Q}
\newcommand{\reals}{\mathbb R}
\newcommand{\complex}{\mathbb C}
\newcommand{\powerset}{\mathscr P}

%-------------------------------------------------------------


\begin{document}
\begin{abstract}
 \end{abstract}

    \title{Reduction of Order}
    
    \maketitle


 Suppose we want to find a general solution to the 2nd order linear differential equation \\ $y''+py'+qy=f$ in the situation where we can find, or are given, a non-trivial solution $y_{1}$ of the complementary homogeneous equation $y''+py'+qy=0$.\\

 Let's look for a solution of the form $y=uy_{1}$. \\

 Then $y'=u'y_{1}+uy_{1}'$ and $y''=u''y_{1}+2u'y_{1}'+uy_{1}''$.  By substitution we get

 \begin{eqnarray*}
  y''+py'+qy=f &\rightarrow & u''y_{1}+2u'y_{1}'+uy_{1}''\,+\,p(u'y_{1}+uy_{1}')\,+\,q(uy_{1})=f \\
   &\rightarrow & u'' y_1 +(2y_1'+py_1)u'+(y_1 ''+py_1 '+qy_1)u=f \\
   &\rightarrow & u'' y_1 +(2y_1'+py_1)u'+(0)u=f \\
   &\rightarrow & u'' y_1 +(2y_1'+py_1)u'=f \\
   &\rightarrow & y_1 z'  +(2y_1'+py_1)z=f\mbox{ using the sub. }z=u'.\\
 \end{eqnarray*}

 This results in a 1st order linear de for $z$, which can be solved for $z=u'$, which can in turn be solved for $u$, and then $y=uy_1$.\\

 This method is called $``$Reduction of Order."  You can choose to memorize \framebox{$y_1 z'  +(2y_1'+py_1)z=f$} but make sure the de is in the standard form!  The process works even when the de isn't in the standard form.  Our examples will illustrate that.\\


 Let's find the general solution to $x^{2}y''+xy'-y=4x^{-2}$ given that $y=x$ is a solution to the complementary equation.\\

 Setting $y=ux$ yields $y'=u'x+u$ and $y''=u''x+2u'$.  By substitution, we get  \\$\displaystyle x^{2}(u''x+2u')+x(u'x+u)-ux=4x^{-2}$ $\mbox{ which implies }x^{3}u''+3x^{2}u'=4x^{-2}$.\\

 In terms of $z=u'$ this is equivalent to $x^3 z'+3x^2 z=4x^{-2}$.\\

 The left side is already in the form $(x^3 z)'$ so $(x^3 z)'=4x^{-2}$ implies $x^3 z=-4x^{-1}+c_1$.

 \newpage
 

 So $z=u'=-4x^{-4}+c_1x^{-3}\rightarrow $ $u=\frac{4}{3}x^{-3}+c_1x^{-2}+c_2$.  (Note: $-\frac{c_1}{2}$ is as arbitrary as $c_1$.) \\ So the general solution is \\$\displaystyle y=ux=\frac{4}{3}x^{-2}+c_1x^{-1}+c_2x=\frac{4}{3x^2}+\frac{c_1}{x}+c_2x=\frac{4+c_1 x+c_2 x^3}{3x^2}$.\\ (Note: $3c_1,3c_2$ are as arbtirary as $c_1,c_2$).



Let's solve the IVP $(x^{2}-4)y''+4xy'+2y=x+2$, $y(0)=0$, $\displaystyle y'(0)=-\frac{3}{8}$ , given that $y=\dfrac{1}{x-2}$ is a complementary solution.  Let $y=u(x-2)^{-1}$. \\

 Then $y'=u'(x-2)^{-1}-u(x-2)^{-2}$ and $y''=u''(x-2)^{-1}-2u'(x-2)^{-2}+2u(x-2)^{-3}$.  So


$$(x^{2}-4)(u''(x-2)^{-1}-2u'(x-2)^{-2}+2u(x-2)^{-3})+4x(u'(x-2)^{-1}-u(x-2)^{-2})+2u(x-2)^{-1}=x+2\rightarrow$$

\medskip

 $$(x+2)u''-2(x+2)(x-2)^{-1}u'+2(x+2)(x-2)^{-2}u+4x(x-2)^{-1}u'-4x(x-2)^{-2}u+2(x-2)^{-1}u=x+2\rightarrow$$


\medskip

 $$(x+2)u''+(-2(x+2)(x-2)^{-1}+4x(x-2)^{-1})u'+(2(x+2)(x-2)^{-2}-4x(x-2)^{-2}+2(x-2)^{-1})u=x+2\rightarrow$$


\medskip

 $(x+2)u''+2u'=x+2$.   Letting $z=u'$ and dividing through by $x+2$ yields $z'+2(x+2)^{-1}z=1$.\\

 We find an integrating factor: $\displaystyle \mu=e^{2\ln{|x+2|}}=(x+2)^{2}$.\\

 Let's solve $(x+2)^{2}z'+2(x+2)z=(x+2)^{2}$, which is equivalent to $((x+2)^{2}z)'=(x+2)^{2}$.\\

 Integrating yields $(x+2)^{2}z=\frac{1}{3}(x+2)^{3}+c_1$.  So $z=u'=\frac{1}{3}(x+2)+c_1(x+2)^{-2}$. \\

 So $u=\frac{1}{6}(x+2)^2+c_1(x+2)^{-1}+c_2$.\\

 Finally we get a general solution by multiplying $u$ by $y_1=(x-2)^{-1}$, so \\$\displaystyle y=\frac{1}{6}(x+2)^2(x-2)^{-1}+c_1(x+2)^{-1}(x-2)^{-1}+c_2(x-2)^{-1}$.  Now we impose the initial conditions:  $y(0)=0$ implies $\displaystyle 0=-\frac{1}{3}-\frac{c_1}{4}-\frac{c_2}{2}$.

 $$y'= \frac{1}{3}(x+2)(x-2)^{-1}-\frac{1}{6}(x+2)^2(x-2)^{-2} -c_1(x+2)^{-2}(x-2)^{-1}-c_1(x+2)^{-1}(x-2)^{-2}-c_2(x-2)^{-2}.$$

 So $\displaystyle y'(0)=-\frac{3}{8}$ implies $\displaystyle -\frac{3}{8}=-\frac{1}{3}-\frac{1}{6} +\frac{c_1}{8}-\frac{c_1}{8}-\frac{c_2}{4}=-\frac{1}{2}-\frac{c_2}{4}$.  So $\displaystyle c_2=-\frac{1}{2}$, and $\displaystyle c_1=-\frac{1}{3}$.\\


 So the IVP solution is $ \displaystyle y=\frac{1}{6}(x+2)^2(x-2)^{-1}-\frac{1}{3}(x+2)^{-1}(x-2)^{-1}-\frac{1}{2}(x-2)^{-1}$,\\

 equivalently, $\displaystyle y=\frac{x(x+3)^2}{6(x^2-4)}$.



\end{document}
